\begin{annexitem}{annex:p3.delta.1}{thm:p03-delta-packing}{Indexed witnesses}
\begin{proof}
We describe the packing explicitly. Let \(D\) be the named reduct
domain in a larger ambient structure. A projective presentation can be
written with a witness domain, private copies of its symbols, and a
unary predicate selecting its reduct. If the original presentation uses
a zero-term interpretation, name its image symbols privately and
require their restrictions to the selected domain to agree with the
intended reduct relations. This is first-order bookkeeping. Function
symbols keep their closure conditions. No relation outside the selected
domain is overwritten.

Suppose the presentation is required for each parameter tuple
\(\bar a\in D^r\) satisfying an activation condition. For each
witness relation of arity \(n\) introduce a private relation of arity
\(n+r\), and for each witness function of arity \(n\) a private
function of arity \(n+r\). This includes constants as arity zero.
Introduce also a private domain relation \(W(x,\bar a)\). A one-step
\(r\)-term interpretation maps the witness universe symbol to \(W\)
and its other symbols to their arity-raised copies, using fresh
parameter constants as its terms. Translate the witness sentence once
by this interpretation. Conjoin the first-order closure guards for
every function and constant of the finite witness signature, including
any that do not occur in the sentence. By the guarded Translation Axiom
and these guards the result asserts exactly that the whole fiber is a
structure and satisfies the original sentence. Universal quantification over the fresh parameter
constants, restricted to the activation condition, requires this at
every active tuple. This uses the individual translation operation and
the quantifier operation, not a composition equation for two
domain-changing symbol maps.

First-order clauses identify the selected reduct of every active fiber
with the required base structure. In the case where the reduct universe
is \(D\), require \(D\subseteq W_{\bar a}\), and require its private
selection predicate, within \(W_{\bar a}\), to hold exactly on \(D\).
Require its private relation symbols, on tuples from \(D\), to agree
with the intended relations at \(\bar a\). For a smaller named reduct
domain use that domain in these clauses instead. Nullary relations
have the same clause with an empty object tuple. Extra values of a
private relation outside its reduct remain unconstrained by these
identifications and can retain everything required by the witness
sentence.

These clauses are sufficient: from a model one reads each actual fiber
and obtains the demanded witness by the Translation Axiom. They are
also necessary in the projective sense. Given the required witnesses,
choose one for each active tuple in the surrounding set theory.
Transport its reduct onto the named domain; tag its other objects by
the tuple and requirement so that they are otherwise disjoint. Take
the union of all these domains, the base domain, and one additional
object. The arity-raised symbols take their witness values on each
active fiber and arbitrary values elsewhere; functions use the
additional object as a default outside their required fibers.
Different tuples have different final argument tuples, so these
prescriptions do not conflict even though all fibers share the base
domain. The extra object also makes the ambient structure nonempty
when the base and every required fiber are empty. There are set many
active tuples, and each chosen witness is a set-sized structure, so
the union is a set. For finitely many requirements use disjoint private
symbols. All of this uses a finite signature and one finite conjunction
of \(\mathcal L\)-sentences.

Apply this construction to a finite generalized-quantifier formula.
For each subformula \(\chi\), name a relation \(T_\chi\) on the
base domain for its free-variable assignments. Impose ordinary
first-order truth clauses for equality, base relation atoms, Boolean
connectives, and first-order quantifiers. For example an existential
node has the clause
\[
 T_{\exists x\chi}(\bar a)\ \longleftrightarrow
 \exists x\in D\,T_\chi(\bar a,x),
\]
with the evident coordinate identifications for repeated or unused
variables. A nullary truth relation records a sentence's truth.

At a node \(Q_K(\chi_1,\ldots,\chi_s)\) and a tuple \(\bar a\)
for its external parameters, the intended \(\rho\)-structure has
universe \(D\) and relations given by the appropriate sections of
\(T_{\chi_i}\). Activate the positive presentation of \(K\) when
this node's truth relation holds at \(\bar a\), and activate its
negative presentation when that relation fails. Pack both requirements
as above. Complementarity of the index and induction up the finite
parse tree force every named truth relation to be its actual semantic
truth relation. Conversely those actual relations, together with
witnesses for the appropriate class at every tuple, satisfy all the
clauses. Finally require the root truth relation, or its negation.
This gives respectively a positive and a negative projective
presentation. Named parameter constants are kept in the selected
reduct and are treated by the same fiber and closure clauses.
\end{proof}
\end{annexitem}

\begin{annexitem}{annex:p3.delta.2}{thm:p03-delta-closure}{Closure and minimality}
\begin{proof}
First note two elementary consequences of the packing construction.
Finite unions and intersections of projective classes are projective:
activate respectively one selected alternative, or every required
presentation, and use private fibers over the same reduct. A
first-order definable class is projective using its own sentence.
Also, successive projective projections can be combined. Keep the
outer and inner witness domains as separate predicates in the final
ambient structure, with private copies of the intermediate signature.
Require the inner reduct domain to lie in the intermediate domain and
its symbols to agree there. Guarded relativization describes each
witness on its own domain. Thus the resulting structure supplies both
successive witnesses. Conversely two witnesses can be transported and
amalgamated over their common intermediate structure. This argument
keeps the two domains explicitly; their naive componentwise symbol
composite is not used.

For a relational \(\mathcal L\)-sentence \(\varphi\), its model
class and its complement have the presentations \(\varphi\) and
\(\neg\varphi\). Their index gives \(Q_K\), applied to the original
atomic relations, equivalent to \(\varphi\). If function symbols
occur, replace them in the input signature by their graphs. The
positive class is: these graphs are total and functional and their
associated functional structure satisfies \(\varphi\). It is
projective by adding the functions, the graph equations, and
\(\varphi\). Its complement is the union of invalid graphs and
valid graphs whose associated structure satisfies \(\neg\varphi\).
This too is projective by the preceding closure under union. Applying
the resulting quantifier to actual function graphs therefore gives an
equivalent sentence in the original signature. This proves extension.

The constructed syntax also has the required individual translations.
For a relational index \(K\), the class of structures with a named
predicate \(D\) whose induced reduct belongs to \(K\) has both
projective presentations: compose the presentation for \(K\), or
its complement, with that selected reduct, retaining the intermediate
domain as above. Its index is therefore available. At a translated
quantifier node use this index to restrict the node's universe to the
selected domain, and translate its relation formulas recursively.
Ordinary first-order quantifiers are restricted to that domain.
Function interpretations add their finite closure guards; the
complement includes the first-order failure of a guard. This gives
exactly the one-step guarded truth condition. Choose the operation
from the symbols actually occurring in the input formula, so
restriction along a source-language inclusion does not change it.
Universe-preserving symbol substitutions are ordinary syntactic
substitutions and compose functorially. Thus the construction stays
within the stated guarded interface, without imposing the invalid
full composition law.

By \cref{thm:p03-delta-packing}, every
\(\Delta(\mathcal L)\)-sentence has complementary projective
presentations in \(\mathcal L\). Conversely a complementary pair on
a finite relational signature is an index, and applying its quantifier
to the signature's own relations defines its positive class. The
function-graph argument handles other finite signatures. This proves
part~3.

A projective class in \(\Delta(\mathcal L)\) has a witnessing
\(\Delta(\mathcal L)\)-sentence. Replace that sentence by its
positive projective presentation in \(\mathcal L\), and combine the
two projections while retaining their domains. The result is a
projective presentation in \(\mathcal L\). The opposite inclusion
follows from extension. This proves part~2.

If a class and its complement are projective in
\(\Delta(\mathcal L)\), part~2 makes both projective in
\(\mathcal L\); part~3 makes the class definable in
\(\Delta(\mathcal L)\). Thus the latter has the Souslin--Kleene
property. Let \(\mathcal M\) be any extension of \(\mathcal L\)
with that property. Each \(\Delta(\mathcal L)\)-definable class
and its complement are projective in \(\mathcal L\), hence
projective in \(\mathcal M\) by replacing their witness sentences
with equivalent \(\mathcal M\)-sentences. Souslin--Kleene then
defines the class in \(\mathcal M\). This proves leastness without
assuming an unproved substitution operation for arbitrary
\(\mathcal M\)-sentences. Monotonicity follows by the same inclusion
of projective presentations; applying leastness to the already closed
extension proves idempotence.
\end{proof}
\end{annexitem}
